\documentclass[letterpaper,10pt]{article}

\usepackage[empty]{fullpage}
\usepackage{titlesec}
\usepackage[dvipsnames]{xcolor}
\usepackage[final]{microtype}
\usepackage[hidelinks]{hyperref}
\usepackage{fancyhdr}
\usepackage{enumitem}
\usepackage[english]{babel}
\usepackage{tabularx}

\pagestyle{fancy}
\fancyhf{}
\renewcommand{\headrulewidth}{0pt}
\setlength{\footskip}{4.08003pt}
\setlength{\parindent}{0pt}

\addtolength{\oddsidemargin}{-0.6in}
\addtolength{\evensidemargin}{-0.6in}
\addtolength{\textwidth}{1.1in}
\addtolength{\topmargin}{-.6in}
\addtolength{\textheight}{1.2in}

% ── SPACING SCALE ───────────────────────────────────────────────────────────
% Change these four numbers to retune the whole document consistently.
% If the resume spills onto a second page, shrink these before cutting content.
\newcommand{\SecBefore}{5pt}    % blank space above a section header
\newcommand{\SecAfter}{4pt}     % space between the header's rule and content
\newcommand{\EntryGap}{4pt}     % space between jobs / degrees
\newcommand{\BulletGap}{2pt}    % space between an entry title and its bullets

% Section header: uppercase, bold, rule below
% NOTE: the rule lives in its own command on purpose. Writing
% \titlerule[0.8pt] directly inside \titleformat's [...] argument breaks --
% TeX reads the ] of [0.8pt] as the end of the argument.
\newcommand{\SecRule}{\titlerule[0.8pt]}
\titleformat{\section}
  {\raggedright\normalsize\bfseries}{}{0em}{\MakeUppercase}
  [\vspace{1pt}\SecRule]
\titlespacing*{\section}{0pt}{\SecBefore}{\SecAfter}

% Two-column row: content left, date/location right
\newcommand{\tworow}[2]{%
  \begin{tabularx}{\textwidth}{@{}>{\raggedright\arraybackslash}X r@{}}
    #1 & #2
  \end{tabularx}\par}

\newcommand{\bulletsStart}{\vspace{\BulletGap}\begin{itemize}[leftmargin=0.2in,noitemsep,topsep=0pt,parsep=0pt]}
\newcommand{\bulletsEnd}{\end{itemize}}
\newcommand{\entrygap}{\vspace{\EntryGap}}

% ── WRITING NOTES ───────────────────────────────────────────────────────────
% - Escape these characters in text: \% \$ \& \# \_
% - Bold the number in each bullet: \textbf{30\%}
% - Approximately: \raise.17ex\hbox{$\scriptstyle\sim$}75 renders as ~75
% - Date ranges use an en dash: 2022 -- 2024
% - Bullet pattern: result first, then how. "Cut X by Y\% by building Z."

\begin{document}

% ══ HEADER ══════════════════════════════════════════════════════════════════
\begin{center}
  {\Large \textbf{Your Name}} \\ \vspace{3pt}
  \small
  City, ST $|$ 555-555-5555 $|$
  \href{mailto:you@example.com}{you@example.com} \\ \vspace{1pt}
  \href{https://www.linkedin.com/in/your-handle/}{linkedin.com/in/your-handle} $|$
  \href{https://github.com/your-handle}{github.com/your-handle} $|$
  \href{https://your-site.com}{your-site.com}
\end{center}

% ══ PROFESSIONAL SUMMARY ════════════════════════════════════════════════════
\section{Professional Summary}
\small
One-line description of who you are and what you build. Two or three sentences backing it
up with your strongest numbers, such as \textbf{N} users served, \textbf{N+} hours saved, or
\textbf{\$N} in revenue. Name the core tools you work in. Close with the background that sets
you apart.

% ══ TECHNICAL KNOWLEDGE ═════════════════════════════════════════════════════
\section{Technical Knowledge}
\small
\begin{tabularx}{\textwidth}{@{}>{\raggedright\arraybackslash}p{1.75in} >{\raggedright\arraybackslash}X@{}}
  \textbf{Languages}:               & Language, Language, Language \\[2pt]
  \textbf{Domain Skills}:           & Skill, Skill, Skill, Skill \\[2pt]
  \textbf{Frameworks \& Libraries}: & Framework, Framework, Framework \\[2pt]
  \textbf{Cloud \& Tools}:          & Tool, Tool, Tool \\[2pt]
  \textbf{Other}:                   & Skill, Skill, Skill \\
\end{tabularx}

% ══ EXPERIENCE ══════════════════════════════════════════════════════════════
\section{Experience}
\small

\tworow{\textbf{Company Name}}{City, ST}
\tworow{Job Title}{\textit{2024 -- Present}}
\bulletsStart
  \item Result with a \textbf{number} by doing the thing, using the tools that matter.
  \item Result with a \textbf{number} by doing the thing, using the tools that matter.
  \item Result with a \textbf{number} by doing the thing, using the tools that matter.
\bulletsEnd

\entrygap

\tworow{\textbf{Company Name}}{City, ST}
\tworow{Job Title}{\textit{2021 -- 2024}}
\bulletsStart
  \item Result with a \textbf{number} by doing the thing, using the tools that matter.
  \item Result with a \textbf{number} by doing the thing, using the tools that matter.
\bulletsEnd

% ══ EDUCATION ═══════════════════════════════════════════════════════════════
\section{Education}
\small

\tworow{\textbf{University Name}, City, ST}{2021}
\tworow{Degree in Field of Study, \textit{Honors if any}}{}

% ══ CERTIFICATIONS & AWARDS ═════════════════════════════════════════════════
% Delete this section if you have nothing to put here.
\section{Certifications \& Awards}
\small
\bulletsStart
  \item \textbf{Award or certification name}, Issuer (Year): one line on what it was for.
\bulletsEnd

\end{document}
